%% USPSC-Apendice.tex
% ---
% Inicia os apêndices
% ---

\begin{apendicesenv}

\chapter{Concurrence of an X-state}
\label{apx:concurrence}

A general way to define an X-state is
\begin{equation}
    \rho = 
    \begin{bmatrix}
        a & 0 & 0 & f\\
        0 & b & e & 0\\
        0 & e^* & c & 0\\
        f^* & 0 & 0 & d
    \end{bmatrix}.
\end{equation}

\noindent As we presented in Sec. \ref{sec:correlations}, the concurrence of two qubits is defined as
\begin{equation}
    C({\rho}) = \max\{ 0, \lambda_1 - \lambda_2 - \lambda_3 - \lambda_4 \},
\end{equation}

\noindent where $\lambda_i > \lambda_j$ for $i<j$ are the square root of the eigenvalues of the matrix $R = \rho \tilde{\rho}$ with $\tilde{\rho} = (\sigma_y \otimes \sigma_y) \rho^* (\sigma_y \otimes \sigma_y)$ and $\rho^*$ being the complex conjugate of the density matrix. To start computing the concurrence of an X-state, we first find that
\begin{equation}
    \rho^* = 
    \begin{bmatrix}
        a & 0 & 0 & f^*\\
        0 & b & e^* & 0\\
        0 & e & c & 0\\
        f & 0 & 0 & d
    \end{bmatrix}.
\end{equation}

\noindent Using that
\begin{equation}
    \sigma_y \otimes \sigma_y = 
    \begin{bmatrix}
        0 & -i\\
        i & 0
    \end{bmatrix} \otimes
    \begin{bmatrix}
        0 & -i\\
        i & 0
    \end{bmatrix} =
    \begin{bmatrix}
        0 & 0 & 0 & -1\\
        0 & 0 & 1 & 0\\
        0 & 1 & 0 & 0\\
        -1 & 0 & 0 & 1
    \end{bmatrix},
\end{equation}

\noindent it is possible to compute
\begin{equation}
    \begin{split}
        \tilde{\rho} &= (\sigma_y \otimes \sigma_y) \rho^* (\sigma_y \otimes \sigma_y)\\
        &= \begin{bmatrix}
        0 & 0 & 0 & -1\\
        0 & 0 & 1 & 0\\
        0 & 1 & 0 & 0\\
        -1 & 0 & 0 & 1
    \end{bmatrix}
    \begin{bmatrix}
        a & 0 & 0 & f^*\\
        0 & b & e^* & 0\\
        0 & e & c & 0\\
        f & 0 & 0 & d
    \end{bmatrix}
    \begin{bmatrix}
        0 & 0 & 0 & -1\\
        0 & 0 & 1 & 0\\
        0 & 1 & 0 & 0\\
        -1 & 0 & 0 & 1
    \end{bmatrix}\\
    &= \begin{bmatrix}
        d & 0 & 0 & f\\
        0 & c & e & 0\\
        0 & e^* & b & 0\\
        f^* & 0 & 0 & a
    \end{bmatrix}.
    \end{split}
\end{equation}

\noindent As a result, we find
\begin{equation}
    \begin{split}
        R &= \rho \tilde{\rho}\\
        &= \begin{bmatrix}
        a & 0 & 0 & f\\
        0 & b & e & 0\\
        0 & e^* & c & 0\\
        f^* & 0 & 0 & d
    \end{bmatrix}
    \begin{bmatrix}
        d & 0 & 0 & f\\
        0 & c & e & 0\\
        0 & e^* & b & 0\\
        f^* & 0 & 0 & a
    \end{bmatrix}\\
    &= \begin{bmatrix}
        ad+|f|^2 & 0 & 0 & 2af\\
        0 & bc+|e|^2 & 2be & 0\\
        0 & 2ce^* & bc+|e|^2 & 0\\
        2df^* & 0 & 0 & ad+|f|^2
    \end{bmatrix}.
    \end{split}
\end{equation}

\noindent To obtain the concurrence, we must use the eigenvalues of $R$: $\lambda_1 = ad+|f|^2 (1 + 2\sqrt{ad})$, $\lambda_2 = ad+|f|^2 (1 - 2\sqrt{ad})$, $\lambda_3 =  bc+|e|^2(1 + 2\sqrt{bc})$, and $\lambda_4 =  bc+|e|^2(1 - 2\sqrt{bc})$. Without specifying the matrix elements of $\rho$, we cannot go further. But the next steps are taking the square root of each eigenvalue of $R$, organizing them in decreasing order and substituting in Eq. \eqref{eq:concurrence}.




%%%%%%%%%%%%%%%%%%%%%%%%%%%%%%%%%%%%%%%%%%%%%%%%%%%%%
\chapter{Derivation of the Wigner-Fisher information}
\label{apx:WFI}

The Wigner-Fisher information can be obtained from similar calculations as the ones of Sec \ref{sec:thermal-kinematics}. First, we consider two close Wigner distributions and make a second order Taylor expansion in a parameter $\theta$ of $W(\mathbf{x}, \theta+d\theta)$, finding that the Wigner relative entropy (defined in the first equality of Eq. \eqref{eq:wigner-relative-entropy}) becomes
\begin{equation}
    \begin{split}
        K[&W(\mathbf{x}, \theta+d\theta) || W(\mathbf{x}, \theta)] = \int W(\mathbf{x}, \theta+d\theta) \ln \left ( \frac{W(\mathbf{x}, \theta+d\theta)}{W(\mathbf{x}, \theta)} \right ) d^2 \mathbf{x} \\
        &= \int [W(\mathbf{x}, \theta) + \partial_\theta W(\mathbf{x}, \theta) d\theta + \frac{1}{2} \partial_\theta^2 W(\mathbf{x}, \theta) d\theta^2 + \mathcal{O}(d\theta^3)]\\
        &\times \left[\frac{\partial_\theta W(\mathbf{x}, \theta)}{W(\mathbf{x}, \theta)} d\theta + \frac{1}{2}\frac{\partial_\theta^2 W(\mathbf{x}, \theta)}{W(\mathbf{x}, \theta)} d\theta^2 - \frac{1}{2} \frac{(\partial_\theta W(\mathbf{x}, \theta)d\theta)^2}{W(\mathbf{x}, \theta)^2} + \mathcal{O}(d\theta^3) \right] d^2 \mathbf{x} \\
        &= \frac{1}{2}\int \frac{(\partial_\theta W(\mathbf{x}, \theta) d\theta)^2}{W(\mathbf{x}, \theta)} d^2 \mathbf{x} + \mathcal{O}(d\theta^3)\\
        &= \frac{1}{2} \mathcal{I}_W(\theta) d\theta^2 + \mathcal{O}(d\theta^3),
    \end{split}
\end{equation}

\noindent where we used $\ln (1+x) = x - \frac{1}{2}x^2 + \mathcal{O}(x^3)$ in the second line and that $\int W(\mathbf{x}, \theta) d^2 \mathbf{x} = 1$ in the third line. In the last equality, we identify the Wigner-Fisher information as 
\begin{equation}
    \mathcal{I}_W(\theta) = \int \frac{(\partial_\theta W(\mathbf{x}, \theta) d\theta)^2}{W(\mathbf{x}, \theta)}.
\end{equation}

\noindent Additionally, we can also write the Wigner-Fisher information in terms of the mean vector and the covariance matrix. One way to do it, is expanding the last line of Eq. \eqref{eq:wigner-relative-entropy}, that in our case is
\begin{equation}
    \begin{split}
        K[W(\mathbf{x}, \theta+d\theta) || W(\mathbf{x}, \theta)] &= -1 + \frac{1}{2} \Big[(\mathbf{v}_{\theta+d\theta} - \mathbf{v}_\theta)^\dagger \Theta_\theta^{-1} (\mathbf{v}_{\theta+d\theta} - \mathbf{v}_\theta) \\
        &+ \text{Tr}(\Theta_\theta^{-1} \Theta_{\theta+d\theta}) + \ln \frac{|\Theta_\theta|}{|\Theta_{\theta+d\theta}|} \Big].
    \end{split}
    \label{eq:K-wigner-fisher}
\end{equation}

\noindent First, we use the second order Taylor expansion of
\begin{align}
    &\mathbf{v}_{\theta+d\theta} \approx \mathbf{v}_{\theta} + \left(\frac{d\mathbf{v}_{\theta}}{d\theta}\right)d\theta + \frac{1}{2}\left(\frac{d^2\mathbf{v}_{\theta}}{d\theta^2}\right)d\theta^2 , \\
    &\Theta_{\theta+d\theta} \approx \Theta_{\theta} + \left(\frac{d\Theta_{\theta}}{d\theta}\right)d\theta + \frac{1}{2}\left(\frac{d^2\Theta_{\theta}}{d\theta^2}\right)d\theta^2 , \\
    &\ln( |\Theta_{\theta+d\theta}|) \approx \ln( |\Theta_{\theta}|) + \left(\frac{d}{d\theta} \ln( |\Theta_{\theta}|)\right)d\theta + \frac{1}{2}\left(\frac{d^2}{d\theta^2} \ln( |\Theta_{\theta}|) \right)d\theta^2,
    \label{eq:expansao-ln-Theta}
\end{align}

\noindent such that the mean vector term in Eq. \eqref{eq:K-wigner-fisher} becomes
\begin{equation}
    (\mathbf{v}_{\theta+d\theta} - \mathbf{v}_\theta)^\dagger \Theta_\theta^{-1} (\mathbf{v}_{\theta+d\theta} - \mathbf{v}_\theta) \approx \dot{\mathbf{v}}_{\theta}^\dagger \Theta_\theta^{-1} \dot{\mathbf{v}}_{\theta} d\theta^2,
\end{equation}

\noindent and the trace term in Eq. \eqref{eq:K-wigner-fisher} becomes
\begin{align}
    \text{Tr}(\Theta_\theta^{-1} \Theta_{\theta+d\theta}) &\approx \text{Tr}\left[\Theta_\theta^{-1} \Theta_{\theta} + \Theta_\theta^{-1} \dot{\Theta}_{\theta} d\theta + \frac{1}{2} \Theta_\theta^{-1} \ddot{\Theta}_{\theta} d\theta^2 \right] \nonumber\\
    &= 2 + \text{Tr}\left[\Theta_\theta^{-1} \dot{\Theta}_{\theta} \right] d\theta + \frac{1}{2} \text{Tr}\left[\Theta_\theta^{-1} \ddot{\Theta}_{\theta} \right] d\theta^2.
\end{align}

\noindent To rewrite the logarithm term in Eq. \eqref{eq:K-wigner-fisher}, we will need the identities \cite{medina-2026}
\begin{align}
    &\frac{d}{d\theta} \ln(|\Theta_\theta|) = \text{Tr}\left[\Theta_\theta^{-1} \dot{\Theta}_{\theta} \right]\\
    &\frac{d^2}{d\theta^2} \ln(|\Theta_\theta|) = \text{Tr}\left[\Theta_\theta^{-1} \ddot{\Theta}_{\theta} + \Theta_\theta^{-1} \dot{\Theta}_{\theta}\Theta_\theta^{-1} \dot{\Theta}_{\theta} \right],
\end{align}

\noindent which can be used in Eq. \eqref{eq:expansao-ln-Theta}, resulting in
\begin{equation}
    \ln( |\Theta_{\theta+d\theta}|) \approx \ln( |\Theta_{\theta}|) + \text{Tr}\left[\Theta_\theta^{-1} \dot{\Theta}_{\theta} \right] d\theta + \frac{1}{2}\text{Tr}\left[\Theta_\theta^{-1} \ddot{\Theta}_{\theta} + \Theta_\theta^{-1} \dot{\Theta}_{\theta}\Theta_\theta^{-1} \dot{\Theta}_{\theta} \right]d\theta^2,
\end{equation}

\noindent that allow us to rewrite the logarithm term as
\begin{equation}
    \ln \frac{|\Theta_\theta|}{|\Theta_{\theta+d\theta}|} \approx -\text{Tr}\left[\Theta_\theta^{-1} \dot{\Theta}_{\theta} \right] d\theta - \frac{1}{2}\text{Tr}\left[\Theta_\theta^{-1} \ddot{\Theta}_{\theta}\right]d\theta^2 + \text{Tr}\left[(\Theta_\theta^{-1} \dot{\Theta}_{\theta})^2 \right]d\theta^2.
\end{equation}

\noindent As a result, the Wigner relative entropy expanded in terms of the mean vector and covariance matrix is
\begin{equation}
    K[W(\mathbf{x}, \theta+d\theta) || W(\mathbf{x}, \theta)] \approx \frac{1}{2} \{\dot{\mathbf{v}}_t^\dagger \Theta^{-1}_t \dot{\mathbf{v}}_t + \text{Tr}[(\Theta_t^{-1} \dot{\Theta}_t)^2]\} d\theta^2,
\end{equation}

\noindent in which we identify the Wigner-Fisher information as
\begin{equation}
    \mathcal{I}_W (t) = \dot{\mathbf{v}}_t^\dagger \Theta^{-1}_t \dot{\mathbf{v}}_t + \text{Tr}[(\Theta_t^{-1} \dot{\Theta}_t)^2].
\end{equation}

\end{apendicesenv}
% ---